\section{What Accuracy Hides}
\label{sec:signal}

Everything in this section is recomputed from the records of
\Cref{sec:results}. No model was called. The records carry the answer key and
the model's answer on every row, so the rate the model chose an option can be
compared against the rate the items called for it.

\subsection{The measure}

Accuracy on a two-option key adds together two different quantities: how well a
model tells the cases apart, and how often it picks one option whatever the
case. A model that answers the first option nine times in ten scores about 0.63
on a set that wants it two times in three, and none of that 0.63 is
discrimination.

An earlier version of this section said ``balanced'' there, which is wrong and
wrong in an instructive direction. Write accuracy as
$\pi \cdot \mathrm{sens} + (1 - \pi) \cdot \mathrm{spec}$ for a base rate $\pi$.
At $\pi = 0.5$ that collapses to $(J + 1)/2$, a function of informedness alone,
so on a balanced key response bias contributes exactly nothing to accuracy and
there is nothing for a decomposition to recover. Every template in this corpus
is balanced by construction, and the identity holds in our records to the fourth
decimal in every cell where no answer went unread. The argument for reporting
$J$ therefore has to be made on other grounds, and the rest of this section
makes it: what accuracy hides here is not bias but which templates carry any
discrimination at all, and what breaks the identity is not bias either but parse
failure that lands unevenly on the key.

Signal detection separates them. \emph{Informedness}, Youden's
$J = \mathrm{sensitivity} + \mathrm{specificity} - 1$, is zero for any
constant-answer policy at any base rate and is unmoved by a pure shift in
preference. It also does not depend on which option is called positive:
swapping the labels exchanges sensitivity with specificity and leaves the sum
alone, so a per-template choice for which no principled basis exists does not
have to be made. \emph{Skew}, the model's rate for an option minus the rate the items
wanted it, names the other half directly.

We compute both per template and average over templates, so no template's item
count decides the answer. Rows the model did not parse are excluded, which is a
limitation we return to. A template whose key holds a single answer class has no
informedness at all, and the implementation refuses rather than returning zero:
recording an undefined cell as zero would put it into a mean as though it were a
measurement of no discrimination.

\subsection{The study, decomposed}

\begin{table}[htbp]
\centering
\signalTable
\caption{Discrimination by arm, with cluster-bootstrap intervals resampling
templates. Parse rate is the share of readings that produced a readable answer;
everything else is computed over that subset.}
\label{tab:signal}
\end{table}

\IfFileExists{figures/signal.tex}{%
\begin{figure}[htbp]
\centering
\section{What Accuracy Hides}
\label{sec:signal}

Everything in this section is recomputed from the records of
\Cref{sec:results}. No model was called. The records carry the answer key and
the model's answer on every row, so the rate the model chose an option can be
compared against the rate the items called for it.

\subsection{The measure}

Accuracy on a two-option key adds together two different quantities: how well a
model tells the cases apart, and how often it picks one option whatever the
case. A model that answers the first option nine times in ten scores about 0.63
on a set that wants it two times in three, and none of that 0.63 is
discrimination.

An earlier version of this section said ``balanced'' there, which is wrong and
wrong in an instructive direction. Write accuracy as
$\pi \cdot \mathrm{sens} + (1 - \pi) \cdot \mathrm{spec}$ for a base rate $\pi$.
At $\pi = 0.5$ that collapses to $(J + 1)/2$, a function of informedness alone,
so on a balanced key response bias contributes exactly nothing to accuracy and
there is nothing for a decomposition to recover. Every template in this corpus
is balanced by construction, and the identity holds in our records to the fourth
decimal in every cell where no answer went unread. The argument for reporting
$J$ therefore has to be made on other grounds, and the rest of this section
makes it: what accuracy hides here is not bias but which templates carry any
discrimination at all, and what breaks the identity is not bias either but parse
failure that lands unevenly on the key.

Signal detection separates them. \emph{Informedness}, Youden's
$J = \mathrm{sensitivity} + \mathrm{specificity} - 1$, is zero for any
constant-answer policy at any base rate and is unmoved by a pure shift in
preference. It also does not depend on which option is called positive:
swapping the labels exchanges sensitivity with specificity and leaves the sum
alone, so a per-template choice for which no principled basis exists does not
have to be made. \emph{Skew}, the model's rate for an option minus the rate the items
wanted it, names the other half directly.

We compute both per template and average over templates, so no template's item
count decides the answer. Rows the model did not parse are excluded, which is a
limitation we return to. A template whose key holds a single answer class has no
informedness at all, and the implementation refuses rather than returning zero:
recording an undefined cell as zero would put it into a mean as though it were a
measurement of no discrimination.

\subsection{The study, decomposed}

\begin{table}[htbp]
\centering
\signalTable
\caption{Discrimination by arm, with cluster-bootstrap intervals resampling
templates. Parse rate is the share of readings that produced a readable answer;
everything else is computed over that subset.}
\label{tab:signal}
\end{table}

\IfFileExists{figures/signal.tex}{%
\begin{figure}[htbp]
\centering
\section{What Accuracy Hides}
\label{sec:signal}

Everything in this section is recomputed from the records of
\Cref{sec:results}. No model was called. The records carry the answer key and
the model's answer on every row, so the rate the model chose an option can be
compared against the rate the items called for it.

\subsection{The measure}

Accuracy on a two-option key adds together two different quantities: how well a
model tells the cases apart, and how often it picks one option whatever the
case. A model that answers the first option nine times in ten scores about 0.63
on a set that wants it two times in three, and none of that 0.63 is
discrimination.

An earlier version of this section said ``balanced'' there, which is wrong and
wrong in an instructive direction. Write accuracy as
$\pi \cdot \mathrm{sens} + (1 - \pi) \cdot \mathrm{spec}$ for a base rate $\pi$.
At $\pi = 0.5$ that collapses to $(J + 1)/2$, a function of informedness alone,
so on a balanced key response bias contributes exactly nothing to accuracy and
there is nothing for a decomposition to recover. Every template in this corpus
is balanced by construction, and the identity holds in our records to the fourth
decimal in every cell where no answer went unread. The argument for reporting
$J$ therefore has to be made on other grounds, and the rest of this section
makes it: what accuracy hides here is not bias but which templates carry any
discrimination at all, and what breaks the identity is not bias either but parse
failure that lands unevenly on the key.

Signal detection separates them. \emph{Informedness}, Youden's
$J = \mathrm{sensitivity} + \mathrm{specificity} - 1$, is zero for any
constant-answer policy at any base rate and is unmoved by a pure shift in
preference. It also does not depend on which option is called positive:
swapping the labels exchanges sensitivity with specificity and leaves the sum
alone, so a per-template choice for which no principled basis exists does not
have to be made. \emph{Skew}, the model's rate for an option minus the rate the items
wanted it, names the other half directly.

We compute both per template and average over templates, so no template's item
count decides the answer. Rows the model did not parse are excluded, which is a
limitation we return to. A template whose key holds a single answer class has no
informedness at all, and the implementation refuses rather than returning zero:
recording an undefined cell as zero would put it into a mean as though it were a
measurement of no discrimination.

\subsection{The study, decomposed}

\begin{table}[htbp]
\centering
\signalTable
\caption{Discrimination by arm, with cluster-bootstrap intervals resampling
templates. Parse rate is the share of readings that produced a readable answer;
everything else is computed over that subset.}
\label{tab:signal}
\end{table}

\IfFileExists{figures/signal.tex}{%
\begin{figure}[htbp]
\centering
\section{What Accuracy Hides}
\label{sec:signal}

Everything in this section is recomputed from the records of
\Cref{sec:results}. No model was called. The records carry the answer key and
the model's answer on every row, so the rate the model chose an option can be
compared against the rate the items called for it.

\subsection{The measure}

Accuracy on a two-option key adds together two different quantities: how well a
model tells the cases apart, and how often it picks one option whatever the
case. A model that answers the first option nine times in ten scores about 0.63
on a set that wants it two times in three, and none of that 0.63 is
discrimination.

An earlier version of this section said ``balanced'' there, which is wrong and
wrong in an instructive direction. Write accuracy as
$\pi \cdot \mathrm{sens} + (1 - \pi) \cdot \mathrm{spec}$ for a base rate $\pi$.
At $\pi = 0.5$ that collapses to $(J + 1)/2$, a function of informedness alone,
so on a balanced key response bias contributes exactly nothing to accuracy and
there is nothing for a decomposition to recover. Every template in this corpus
is balanced by construction, and the identity holds in our records to the fourth
decimal in every cell where no answer went unread. The argument for reporting
$J$ therefore has to be made on other grounds, and the rest of this section
makes it: what accuracy hides here is not bias but which templates carry any
discrimination at all, and what breaks the identity is not bias either but parse
failure that lands unevenly on the key.

Signal detection separates them. \emph{Informedness}, Youden's
$J = \mathrm{sensitivity} + \mathrm{specificity} - 1$, is zero for any
constant-answer policy at any base rate and is unmoved by a pure shift in
preference. It also does not depend on which option is called positive:
swapping the labels exchanges sensitivity with specificity and leaves the sum
alone, so a per-template choice for which no principled basis exists does not
have to be made. \emph{Skew}, the model's rate for an option minus the rate the items
wanted it, names the other half directly.

We compute both per template and average over templates, so no template's item
count decides the answer. Rows the model did not parse are excluded, which is a
limitation we return to. A template whose key holds a single answer class has no
informedness at all, and the implementation refuses rather than returning zero:
recording an undefined cell as zero would put it into a mean as though it were a
measurement of no discrimination.

\subsection{The study, decomposed}

\begin{table}[htbp]
\centering
\signalTable
\caption{Discrimination by arm, with cluster-bootstrap intervals resampling
templates. Parse rate is the share of readings that produced a readable answer;
everything else is computed over that subset.}
\label{tab:signal}
\end{table}

\IfFileExists{figures/signal.tex}{%
\begin{figure}[htbp]
\centering
\input{figures/signal}
\caption{Discrimination against an empty prompt, with cluster-bootstrap
intervals over templates. Every interval crosses zero. Ten clusters is few, and
the widths say so.}
\label{fig:signal}
\end{figure}}{}

\textbf{The study's conclusion survives the correction and gets slightly
stronger.} Every arm's discrimination advantage over an empty prompt has an
interval containing zero, exactly as every accuracy comparison did after Holm.
SkillOpt keeps first place on both measures, at
$\Delta J = \NUM{\deltaJSkillopt}$ with an interval of
$[\NUM{\deltaJSkilloptLow}, \NUM{\deltaJSkilloptHigh}]$ and the widest interval
of any arm here. It is still indistinguishable from nothing.

Two things this does not say. Mean skew runs \NUM{\meanSkewGepa} to
\NUM{\meanSkewOn} in \emph{every} arm including the empty prompt, so the
response bias here is a property of the model on this corpus rather than
something an arm introduced. And $J$ is computed over parsed rows, so
\texttt{gepa}, which parses least at \NUM{\parseRateGepa}, is measured on the
subset it managed to answer, and \Cref{sec:parsefail} shows that subset is not a
random half of its items.

\subsection{Three of ten templates are not measuring anything}

\begin{table}[htbp]
\centering
\templateTable
\caption{Informedness per template, every arm. Templates marked $\dagger$ sit
below $J = \NUM{\lowSignalThreshold}$ in all five arms.}
\label{tab:templates}
\end{table}

\NUM{\lowSignalTemplates} of \NUM{\studyTemplates} templates sit below
$J = \NUM{\lowSignalThreshold}$ in all five arms, and two of them are
indistinguishable from zero. On those items the model is not deciding. What
reaches the headline accuracy is the base rate and whichever way the model
happens to lean.

They hold \NUM{\lowSignalItems} of the \NUM{\totalItems} items, so the run
carries the discriminative signal of the other \NUM{\signalItems}. The power
calculation behind the study assumed \NUM{\totalItems}. That does not change a
result that failed to reject, since less power only makes rejection harder, but
a study designed off those numbers would be sized wrong.

One of the three is the template an earlier entry in our own notebook had called
the hardest in the corpus, on the strength of its low accuracy. It is not hard.
It is unanswered, and accuracy cannot tell those apart.

\subsection{Why we think this generalises}

The two-option forced choice is the workhorse format of this literature, because
it makes ground truth cheap and scoring unambiguous. It is also the format in
which a template that measures nothing is easiest to leave in, for two reasons
that pull in different directions.

The first is aggregation. On a balanced key accuracy over the parsed rows is
$(J+1)/2$, so a template at $J = 0$ reads exactly 0.5000 and is obvious the
moment anyone looks at it per template. Almost nobody does. Pooled over the
\NUM{\templatesSeen} seen templates, two of which sit below
$J = \NUM{\lowSignalThreshold}$ in every arm and one of those at chance, the
same corpus returns \NUM{\accSeenOff} for an empty prompt, which looks like a
corpus that works.

The second is that the metric this study reports is not the parsed one. Counting
an unreadable answer as wrong breaks the identity, and it breaks the ordering
with it. \texttt{rel-003-oncall-escalate} discriminates perfectly for
\texttt{off} in \Cref{tab:templates}, on the
\NUM{\parseRateSeenRelZeroZeroThreeOncallEscalateOff} of items it answered, and
reports \NUM{\accSeenRelZeroZeroThreeOncallEscalateOff}.
\texttt{rel-006-refund-request} barely discriminates at all and reports
\NUM{\accSeenRelZeroZeroSixRefundRequestOff}, higher. On the published number a
template that is always right when it answers ranks below one that is mostly
guessing.

An unbalanced key adds the bias channel on top, and then an arm that shifts a
model's preference toward the option the items happen to favour gains accuracy
without deciding any better. Our own corpus is balanced, so that channel is
closed here and we claim nothing from it. What is open here, and generic, is the
third channel: an arm whose unreadable answers correlate with the key moves
accuracy on the parsed subset by \NUM{\gapParsedUnseenRelZeroZeroOneVendorOutageGepa}
on one template while moving the reported figure the other way
(\Cref{sec:parsefail}). The decomposition costs nothing. It runs over records a completed study already
has, needs no additional calls, and would have taken us from ``ten templates''
to ``seven templates and three that measure nothing'' before we sized anything.

\subsection{What it changed about our own practice}

The repository's corpus-admission rule was accuracy on the control arm inside
$[0.35, 0.75]$. That criterion admits a one-sided template, because
one-sidedness lands squarely inside the band. It has been extended with two
more: skew near zero, so the model is answering from the item, and informedness
below one, so there is something left for an arm to gain. Both are computed from
the calibration run the first criterion already requires. Neither closes the
case, and \Cref{sec:ceiling} gives the template of ours that passes all three
and still measures nothing.

\caption{Discrimination against an empty prompt, with cluster-bootstrap
intervals over templates. Every interval crosses zero. Ten clusters is few, and
the widths say so.}
\label{fig:signal}
\end{figure}}{}

\textbf{The study's conclusion survives the correction and gets slightly
stronger.} Every arm's discrimination advantage over an empty prompt has an
interval containing zero, exactly as every accuracy comparison did after Holm.
SkillOpt keeps first place on both measures, at
$\Delta J = \NUM{\deltaJSkillopt}$ with an interval of
$[\NUM{\deltaJSkilloptLow}, \NUM{\deltaJSkilloptHigh}]$ and the widest interval
of any arm here. It is still indistinguishable from nothing.

Two things this does not say. Mean skew runs \NUM{\meanSkewGepa} to
\NUM{\meanSkewOn} in \emph{every} arm including the empty prompt, so the
response bias here is a property of the model on this corpus rather than
something an arm introduced. And $J$ is computed over parsed rows, so
\texttt{gepa}, which parses least at \NUM{\parseRateGepa}, is measured on the
subset it managed to answer, and \Cref{sec:parsefail} shows that subset is not a
random half of its items.

\subsection{Three of ten templates are not measuring anything}

\begin{table}[htbp]
\centering
\templateTable
\caption{Informedness per template, every arm. Templates marked $\dagger$ sit
below $J = \NUM{\lowSignalThreshold}$ in all five arms.}
\label{tab:templates}
\end{table}

\NUM{\lowSignalTemplates} of \NUM{\studyTemplates} templates sit below
$J = \NUM{\lowSignalThreshold}$ in all five arms, and two of them are
indistinguishable from zero. On those items the model is not deciding. What
reaches the headline accuracy is the base rate and whichever way the model
happens to lean.

They hold \NUM{\lowSignalItems} of the \NUM{\totalItems} items, so the run
carries the discriminative signal of the other \NUM{\signalItems}. The power
calculation behind the study assumed \NUM{\totalItems}. That does not change a
result that failed to reject, since less power only makes rejection harder, but
a study designed off those numbers would be sized wrong.

One of the three is the template an earlier entry in our own notebook had called
the hardest in the corpus, on the strength of its low accuracy. It is not hard.
It is unanswered, and accuracy cannot tell those apart.

\subsection{Why we think this generalises}

The two-option forced choice is the workhorse format of this literature, because
it makes ground truth cheap and scoring unambiguous. It is also the format in
which a template that measures nothing is easiest to leave in, for two reasons
that pull in different directions.

The first is aggregation. On a balanced key accuracy over the parsed rows is
$(J+1)/2$, so a template at $J = 0$ reads exactly 0.5000 and is obvious the
moment anyone looks at it per template. Almost nobody does. Pooled over the
\NUM{\templatesSeen} seen templates, two of which sit below
$J = \NUM{\lowSignalThreshold}$ in every arm and one of those at chance, the
same corpus returns \NUM{\accSeenOff} for an empty prompt, which looks like a
corpus that works.

The second is that the metric this study reports is not the parsed one. Counting
an unreadable answer as wrong breaks the identity, and it breaks the ordering
with it. \texttt{rel-003-oncall-escalate} discriminates perfectly for
\texttt{off} in \Cref{tab:templates}, on the
\NUM{\parseRateSeenRelZeroZeroThreeOncallEscalateOff} of items it answered, and
reports \NUM{\accSeenRelZeroZeroThreeOncallEscalateOff}.
\texttt{rel-006-refund-request} barely discriminates at all and reports
\NUM{\accSeenRelZeroZeroSixRefundRequestOff}, higher. On the published number a
template that is always right when it answers ranks below one that is mostly
guessing.

An unbalanced key adds the bias channel on top, and then an arm that shifts a
model's preference toward the option the items happen to favour gains accuracy
without deciding any better. Our own corpus is balanced, so that channel is
closed here and we claim nothing from it. What is open here, and generic, is the
third channel: an arm whose unreadable answers correlate with the key moves
accuracy on the parsed subset by \NUM{\gapParsedUnseenRelZeroZeroOneVendorOutageGepa}
on one template while moving the reported figure the other way
(\Cref{sec:parsefail}). The decomposition costs nothing. It runs over records a completed study already
has, needs no additional calls, and would have taken us from ``ten templates''
to ``seven templates and three that measure nothing'' before we sized anything.

\subsection{What it changed about our own practice}

The repository's corpus-admission rule was accuracy on the control arm inside
$[0.35, 0.75]$. That criterion admits a one-sided template, because
one-sidedness lands squarely inside the band. It has been extended with two
more: skew near zero, so the model is answering from the item, and informedness
below one, so there is something left for an arm to gain. Both are computed from
the calibration run the first criterion already requires. Neither closes the
case, and \Cref{sec:ceiling} gives the template of ours that passes all three
and still measures nothing.

\caption{Discrimination against an empty prompt, with cluster-bootstrap
intervals over templates. Every interval crosses zero. Ten clusters is few, and
the widths say so.}
\label{fig:signal}
\end{figure}}{}

\textbf{The study's conclusion survives the correction and gets slightly
stronger.} Every arm's discrimination advantage over an empty prompt has an
interval containing zero, exactly as every accuracy comparison did after Holm.
SkillOpt keeps first place on both measures, at
$\Delta J = \NUM{\deltaJSkillopt}$ with an interval of
$[\NUM{\deltaJSkilloptLow}, \NUM{\deltaJSkilloptHigh}]$ and the widest interval
of any arm here. It is still indistinguishable from nothing.

Two things this does not say. Mean skew runs \NUM{\meanSkewGepa} to
\NUM{\meanSkewOn} in \emph{every} arm including the empty prompt, so the
response bias here is a property of the model on this corpus rather than
something an arm introduced. And $J$ is computed over parsed rows, so
\texttt{gepa}, which parses least at \NUM{\parseRateGepa}, is measured on the
subset it managed to answer, and \Cref{sec:parsefail} shows that subset is not a
random half of its items.

\subsection{Three of ten templates are not measuring anything}

\begin{table}[htbp]
\centering
\templateTable
\caption{Informedness per template, every arm. Templates marked $\dagger$ sit
below $J = \NUM{\lowSignalThreshold}$ in all five arms.}
\label{tab:templates}
\end{table}

\NUM{\lowSignalTemplates} of \NUM{\studyTemplates} templates sit below
$J = \NUM{\lowSignalThreshold}$ in all five arms, and two of them are
indistinguishable from zero. On those items the model is not deciding. What
reaches the headline accuracy is the base rate and whichever way the model
happens to lean.

They hold \NUM{\lowSignalItems} of the \NUM{\totalItems} items, so the run
carries the discriminative signal of the other \NUM{\signalItems}. The power
calculation behind the study assumed \NUM{\totalItems}. That does not change a
result that failed to reject, since less power only makes rejection harder, but
a study designed off those numbers would be sized wrong.

One of the three is the template an earlier entry in our own notebook had called
the hardest in the corpus, on the strength of its low accuracy. It is not hard.
It is unanswered, and accuracy cannot tell those apart.

\subsection{Why we think this generalises}

The two-option forced choice is the workhorse format of this literature, because
it makes ground truth cheap and scoring unambiguous. It is also the format in
which a template that measures nothing is easiest to leave in, for two reasons
that pull in different directions.

The first is aggregation. On a balanced key accuracy over the parsed rows is
$(J+1)/2$, so a template at $J = 0$ reads exactly 0.5000 and is obvious the
moment anyone looks at it per template. Almost nobody does. Pooled over the
\NUM{\templatesSeen} seen templates, two of which sit below
$J = \NUM{\lowSignalThreshold}$ in every arm and one of those at chance, the
same corpus returns \NUM{\accSeenOff} for an empty prompt, which looks like a
corpus that works.

The second is that the metric this study reports is not the parsed one. Counting
an unreadable answer as wrong breaks the identity, and it breaks the ordering
with it. \texttt{rel-003-oncall-escalate} discriminates perfectly for
\texttt{off} in \Cref{tab:templates}, on the
\NUM{\parseRateSeenRelZeroZeroThreeOncallEscalateOff} of items it answered, and
reports \NUM{\accSeenRelZeroZeroThreeOncallEscalateOff}.
\texttt{rel-006-refund-request} barely discriminates at all and reports
\NUM{\accSeenRelZeroZeroSixRefundRequestOff}, higher. On the published number a
template that is always right when it answers ranks below one that is mostly
guessing.

An unbalanced key adds the bias channel on top, and then an arm that shifts a
model's preference toward the option the items happen to favour gains accuracy
without deciding any better. Our own corpus is balanced, so that channel is
closed here and we claim nothing from it. What is open here, and generic, is the
third channel: an arm whose unreadable answers correlate with the key moves
accuracy on the parsed subset by \NUM{\gapParsedUnseenRelZeroZeroOneVendorOutageGepa}
on one template while moving the reported figure the other way
(\Cref{sec:parsefail}). The decomposition costs nothing. It runs over records a completed study already
has, needs no additional calls, and would have taken us from ``ten templates''
to ``seven templates and three that measure nothing'' before we sized anything.

\subsection{What it changed about our own practice}

The repository's corpus-admission rule was accuracy on the control arm inside
$[0.35, 0.75]$. That criterion admits a one-sided template, because
one-sidedness lands squarely inside the band. It has been extended with two
more: skew near zero, so the model is answering from the item, and informedness
below one, so there is something left for an arm to gain. Both are computed from
the calibration run the first criterion already requires. Neither closes the
case, and \Cref{sec:ceiling} gives the template of ours that passes all three
and still measures nothing.

\caption{Discrimination against an empty prompt, with cluster-bootstrap
intervals over templates. Every interval crosses zero. Ten clusters is few, and
the widths say so.}
\label{fig:signal}
\end{figure}}{}

\textbf{The study's conclusion survives the correction and gets slightly
stronger.} Every arm's discrimination advantage over an empty prompt has an
interval containing zero, exactly as every accuracy comparison did after Holm.
SkillOpt keeps first place on both measures, at
$\Delta J = \NUM{\deltaJSkillopt}$ with an interval of
$[\NUM{\deltaJSkilloptLow}, \NUM{\deltaJSkilloptHigh}]$ and the widest interval
of any arm here. It is still indistinguishable from nothing.

Two things this does not say. Mean skew runs \NUM{\meanSkewGepa} to
\NUM{\meanSkewOn} in \emph{every} arm including the empty prompt, so the
response bias here is a property of the model on this corpus rather than
something an arm introduced. And $J$ is computed over parsed rows, so
\texttt{gepa}, which parses least at \NUM{\parseRateGepa}, is measured on the
subset it managed to answer, and \Cref{sec:parsefail} shows that subset is not a
random half of its items.

\subsection{Three of ten templates are not measuring anything}

\begin{table}[htbp]
\centering
\templateTable
\caption{Informedness per template, every arm. Templates marked $\dagger$ sit
below $J = \NUM{\lowSignalThreshold}$ in all five arms.}
\label{tab:templates}
\end{table}

\NUM{\lowSignalTemplates} of \NUM{\studyTemplates} templates sit below
$J = \NUM{\lowSignalThreshold}$ in all five arms, and two of them are
indistinguishable from zero. On those items the model is not deciding. What
reaches the headline accuracy is the base rate and whichever way the model
happens to lean.

They hold \NUM{\lowSignalItems} of the \NUM{\totalItems} items, so the run
carries the discriminative signal of the other \NUM{\signalItems}. The power
calculation behind the study assumed \NUM{\totalItems}. That does not change a
result that failed to reject, since less power only makes rejection harder, but
a study designed off those numbers would be sized wrong.

One of the three is the template an earlier entry in our own notebook had called
the hardest in the corpus, on the strength of its low accuracy. It is not hard.
It is unanswered, and accuracy cannot tell those apart.

\subsection{Why we think this generalises}

The two-option forced choice is the workhorse format of this literature, because
it makes ground truth cheap and scoring unambiguous. It is also the format in
which a template that measures nothing is easiest to leave in, for two reasons
that pull in different directions.

The first is aggregation. On a balanced key accuracy over the parsed rows is
$(J+1)/2$, so a template at $J = 0$ reads exactly 0.5000 and is obvious the
moment anyone looks at it per template. Almost nobody does. Pooled over the
\NUM{\templatesSeen} seen templates, two of which sit below
$J = \NUM{\lowSignalThreshold}$ in every arm and one of those at chance, the
same corpus returns \NUM{\accSeenOff} for an empty prompt, which looks like a
corpus that works.

The second is that the metric this study reports is not the parsed one. Counting
an unreadable answer as wrong breaks the identity, and it breaks the ordering
with it. \texttt{rel-003-oncall-escalate} discriminates perfectly for
\texttt{off} in \Cref{tab:templates}, on the
\NUM{\parseRateSeenRelZeroZeroThreeOncallEscalateOff} of items it answered, and
reports \NUM{\accSeenRelZeroZeroThreeOncallEscalateOff}.
\texttt{rel-006-refund-request} barely discriminates at all and reports
\NUM{\accSeenRelZeroZeroSixRefundRequestOff}, higher. On the published number a
template that is always right when it answers ranks below one that is mostly
guessing.

An unbalanced key adds the bias channel on top, and then an arm that shifts a
model's preference toward the option the items happen to favour gains accuracy
without deciding any better. Our own corpus is balanced, so that channel is
closed here and we claim nothing from it. What is open here, and generic, is the
third channel: an arm whose unreadable answers correlate with the key moves
accuracy on the parsed subset by \NUM{\gapParsedUnseenRelZeroZeroOneVendorOutageGepa}
on one template while moving the reported figure the other way
(\Cref{sec:parsefail}). The decomposition costs nothing. It runs over records a completed study already
has, needs no additional calls, and would have taken us from ``ten templates''
to ``seven templates and three that measure nothing'' before we sized anything.

\subsection{What it changed about our own practice}

The repository's corpus-admission rule was accuracy on the control arm inside
$[0.35, 0.75]$. That criterion admits a one-sided template, because
one-sidedness lands squarely inside the band. It has been extended with two
more: skew near zero, so the model is answering from the item, and informedness
below one, so there is something left for an arm to gain. Both are computed from
the calibration run the first criterion already requires. Neither closes the
case, and \Cref{sec:ceiling} gives the template of ours that passes all three
and still measures nothing.
